\chapter*{历史注记}
\phantomsection
\addcontentsline{toc}{chapter}{历史注记}

（注意 --- 罗马数字指本注记末尾的参考文献。）

两个正量的算术平均与几何平均的概念可追溯到古代，特别是毕达哥拉斯学派曾把它们作为其喜爱的论题之一。这两个平均之间的不等式 \(\sqrt{ab} \leqslant \frac{1}{2}(a + b)\) 很可能也为他们所熟知；无论如何，Euclid 在与两数之和给定时使其乘积最大的问题相联系时证明了它。看来人们要等到相当晚的 1729 年，才由 MacLaurin 明确陈述了这一问题向 n 个数的推广以及相应的不等式（尽管比这困难得多的极值问题早在很久以前就已被研究过）。

从 18 世纪末起，其他一些类似的不等式出现在分析与几何中。例如，1789 年，Lhuilier 解决了与不等式

\[
\sqrt {\left(\sum_ {k = 1} ^ {n} a _ {k}\right) ^ {2} + \left(\sum_ {k = 1} ^ {n} b _ {k}\right) ^ {2}} \leqslant \sum_ {k = 1} ^ {n} \sqrt {a _ {k} ^ {2} + b _ {k} ^ {2}};
\]

相联系的一个几何最大值问题；Cauchy 在 1821 年证明了

\[
\left(\sum_ {k = 1} ^ {n} a _ {k} b _ {k}\right) ^ {2} \leqslant \left(\sum_ {k = 1} ^ {n} a _ {k} ^ {2}\right) \left(\sum_ {k = 1} ^ {n} b _ {k} ^ {2}\right)
\]

Hölder 不等式的特殊情形；Buniakowsky 于 1859 年、H. A. Schwarz 于 1885 年把 Cauchy 不等式推广到了积分。

然而，直到 19 世纪末，凸性不等式才成为系统研究的对象。O. Hölder 于 1888 年 (I) 在这个问题中引入了单变量凸函数的思想：他从 \(\log x\) 的凹性出发得到了几何平均不等式；把同样的思想应用于函数 \(x^{r}\) ，他证明了以他的名字命名的不等式，而 L. J. Rogers 在一年之前已经从几何平均不等式出发得到了这一不等式（参见习题 4）。至于 Minkowski 不等式，则由后者于 1896 年（对有限和）在与他关于“数的几何”的著名工作 ((II), pp. 115-117) 的联系中给出证明；虽然凸性思想（对任意多个变量的函数而言）是这项工作的基本概念之一，令人相当奇怪的是，Minkowski 竟未注意到他的不等式可以由 Hölder 方法应用于 \((1 + x^{r})^{1 / r}\) 而得到（他只限于借助无穷小分析寻求极值的经典方法）。

要想更深入地研究凸性不等式及其应用，读者可查阅 Hardy、Littlewood 与 Pólya 专论这一问题的书，其中还含有非常完全的参考文献 (III)。

\begin{enumerate}
\def\labelenumi{(\Roman{enumi})}
\item
  O. HÖLDER, Ueber einen Mittelwertsatz, Göttinger Nachrichten (1889), pp. 38-47.
\item
  H. MINKOWSKI, Geometrie der Zahlen, 第 2 版, Leipzig--Berlin (Teubner), 1910.
\item
  G. H. HARDY, J. E. LITTLEWOOD, G. PÓLYA, Inequalities, Cambridge (University Press), 1934.
\end{enumerate}

